\documentclass[12pt]{article}
\usepackage[margin=1in]{geometry}
\usepackage{booktabs}
\usepackage[round]{natbib}
\usepackage{setspace}
\setstretch{1.35}
\usepackage[hidelinks]{hyperref}

\title{Two Configurations for Reliable Handoffs: A Crisp-Set Comparison of Sixteen Hospital Units\\[1em]
\normalsize\textbf{SYNTHETIC TEST PAPER, NOT REAL RESEARCH}}
\author{Ada Quillfeather and Ramon Teskey\\
\small Department of Organization Studies, Northvale University (a fictional institution)}
\date{}

\begin{document}
\maketitle
\thispagestyle{empty}

\noindent\textit{Note.} This paper is a synthetic test document written to check a review tool. The hospital, its units, the authors, and every informant are fictional, and the case data are synthetic, coded from fictional interviews. No real organization or person is described.

\begin{abstract}
\noindent Shift-change handoffs are a frequent source of error in hospital care, yet units that face similar demands differ widely in how reliably they hand patients over. We studied all 16 inpatient units of Brenmoor Regional Hospital, a fictional teaching hospital, through 74 interviews with nurses, physicians, and nurse managers. We coded each unit into crisp sets from the interview evidence and compared the units with crisp-set qualitative comparative analysis. Eight units sustained reliable handoffs. Two configurations are sufficient for reliable handoffs. We call them the protocol-led unit and the champion-led unit. In the first, a written protocol that staff follow at every shift change carries the handoff without a named champion. In the second, a champion carries it without a protocol. Together the two configurations cover two of the eight reliable units. In this hospital, a protocol without a champion and a champion without a protocol are each part of a configuration sufficient for reliable handoffs.
\end{abstract}

\newpage
\section{Introduction}

A handoff moves responsibility for a patient from one clinician to the next. When it fails, the next shift may not learn of a pending test, a changed medication, or a worsening sign. Hospitals have answered with standard forms, electronic tools, and training. Yet units in the same hospital, with the same forms and the same tools, still differ in whether handoffs hold up week after week.

Two explanations recur in the conversations we had with clinicians. One locates reliability in rules: a written protocol, followed at every change of shift, makes the handoff independent of who happens to be working. The other locates it in people: a respected clinician who watches handoffs and corrects lapses keeps the practice alive. These explanations are usually treated as rivals, or as ingredients that add up. We ask a different question. Which combinations of unit conditions are sufficient for reliable handoffs, and does a unit need both a rule and a person?

Combinations call for a method built for them. Qualitative comparative analysis treats each case as a configuration of conditions and asks which configurations are subsets of the outcome \citep{ragin1987,ragin2008}. It suits a study of a modest number of units known in depth, where the researcher can code each unit from close knowledge of it \citep{fiss2011,greckhamer2018}. We studied every inpatient unit of one fictional hospital, coded each unit from interviews, and compared them.

\section{Rules, people, and the reliable handoff}

Studies of high-reliability work describe two ways that a practice becomes routine. Formal rules write the practice down and make it the same for every shift. They work when staff use them as written and lose force when the form lives in a binder. Informal leadership works through a person whom others watch and copy. It works while that person is present and respected and fades when the person leaves.

The two need not combine. A unit with a protocol that every nurse follows may not need a champion, because the protocol does the work that a champion would otherwise do. A unit without a protocol may still hand over well if a champion keeps the practice on track. If this reasoning holds, configurations sufficient for reliable handoffs should contain a protocol or a champion, and none should need both.

Other conditions may matter in combination with these two. Daily safety huddles give staff a place to raise problems before the shift begins. Adequate staffing leaves time for a full handoff. Long tenure means nurses know the unit's patients and each other. An electronic handoff tool records what was said. A nurse manager present at shift change can see lapses as they happen. We include these five conditions because informants raised each of them, unprompted, as something that helped or hurt handoffs on their unit.

\section{Methods}

\subsection{Setting and data}

Brenmoor Regional Hospital is a fictional 600-bed teaching hospital. We studied all 16 of its inpatient units, so the selection of units follows from the hospital's boundary and no unit was chosen for its outcome. The units include medical and surgical wards and two step-down units. All 16 use the same paper handoff form, and 9 of them also record handoffs in the hospital's electronic tool.

Over one year we conducted 74 semi-structured interviews: 39 with staff nurses, 19 with physicians, and 16 with nurse managers, at least four interviews per unit. We asked informants to describe recent handoffs, what went well, and what went wrong. We recorded and transcribed every interview. Excerpts below are attributed to the informant's role and an interview number.

\subsection{Analytic approach}

We analyze 16 units on seven conditions. The outcome is reliable handoffs. The conditions are a written protocol in use, a handoff champion, daily huddles, adequate staffing, long nurse tenure, an electronic handoff tool, and a nurse manager present at shift change. We first tested each condition and its negation for necessity, as a separate step \citep{schneider2012}, and then analyzed sufficiency. For sufficiency we coded a row as sufficient for the outcome when its raw consistency was at least 0.80, with a frequency cutoff of one unit, which keeps all 16 units. Analyses were run in fsQCA software. We report the conservative solution, which makes no simplifying assumptions about combinations of conditions that no unit shows, because we wanted every configuration we report to rest on units we studied. All 16 units enter the analysis, and every unit has a membership of 0 or 1 in every set.

\subsection{Calibration of the sets}

All sets are crisp. We coded each unit into each set case by case, from the interview evidence, using a written coding rule for each set. Each rule below states the evidence behind it. We drew the line for each rule where the interview evidence divided the units, and we say for each set why the line falls where it does.

\textit{Reliable handoffs (outcome).} A unit is in the set when informants in at least two roles, nurses and physicians, described handoffs over the study year as complete and routine, and no informant described a handoff gap that reached a patient. We required two roles because, in every unit where only one role described handoffs as complete, informants in the other role described recurring gaps. A physician put the test this way: ``I know handoffs work on a unit when I stop having to call back at seven to ask what happened overnight'' (physician, interview 22).

\textit{Protocol.} A unit is in the set when informants described a written handoff protocol that staff followed at every shift change. A protocol kept only on file does not count. We set the rule at use, not existence, because every unit had a protocol on paper, and informants drew the distinction themselves: ``We have a protocol. It's in the blue binder. Nobody opens the blue binder'' (staff nurse, interview 8).

\textit{Champion.} A unit is in the set when informants in at least two roles named the same clinician as the person who kept handoffs on track. We required agreement across roles because a single nurse's praise of a colleague did not, in our interviews, coincide with that colleague correcting other people's handoffs.

\textit{Huddles.} A unit is in the set when informants described a safety huddle held at the start of most shifts. Units where huddles were held only during inspections or after an incident are out of the set, because informants in those units described the huddle as a response to trouble, not a routine.

\textit{Staffing.} A unit is in the set when informants in at least two roles described staffing as adequate on most shifts and no informant described routine borrowing of nurses from other units. Informants treated borrowing as the sign that staffing had broken down, so we used it as the dividing line.

\textit{Tenure.} A unit is in the set when most of the nurses we interviewed had more than three years on the unit. Nurses described the first three years as the period in which they still learned a unit's patients and physicians, which is why the line falls at three years.

\textit{Digital.} A unit is in the set when handoffs were recorded in the electronic handoff tool. Units that handed over on paper or orally only are out of the set. In all but one unit, informants described the tool as either used at every handoff or not used at all, so the rule separated those units cleanly.

\textit{Leader.} A unit is in the set when informants described the nurse manager as present at shift change on most days. We used presence at shift change, not general visibility, because informants who praised a visible manager still described handoffs the manager never saw.

Table~\ref{tab:cases} prints each unit's memberships. Eight of the 16 units are in the outcome set, and each condition holds for between 5 and 10 of the 16 units. Two codes were borderline in our memos: the electronic tool in unit U07, which handed over partly on paper, and nurse tenure in unit U02. Both units are outside the outcome set.

\begin{table}[htbp]
\centering
\caption{Crisp-set memberships of the 16 units (1 = in the set, 0 = out of the set)}
\label{tab:cases}
\begin{tabular}{lcccccccc}
\toprule
Unit & Reliable & Protocol & Champion & Huddles & Staffing & Tenure & Digital & Leader \\
\midrule
U01 & 0 & 1 & 1 & 0 & 1 & 1 & 1 & 0 \\
U02 & 0 & 0 & 0 & 0 & 0 & 1 & 0 & 1 \\
U03 & 1 & 0 & 0 & 1 & 1 & 1 & 1 & 0 \\
U04 & 1 & 0 & 1 & 1 & 1 & 0 & 0 & 0 \\
U05 & 1 & 1 & 1 & 0 & 1 & 1 & 1 & 0 \\
U06 & 1 & 0 & 1 & 0 & 1 & 0 & 1 & 0 \\
U07 & 0 & 1 & 1 & 0 & 1 & 1 & 0 & 0 \\
U08 & 1 & 1 & 0 & 1 & 0 & 0 & 1 & 0 \\
U09 & 1 & 1 & 0 & 1 & 0 & 1 & 0 & 1 \\
U10 & 1 & 1 & 0 & 0 & 0 & 1 & 0 & 1 \\
U11 & 0 & 0 & 0 & 1 & 0 & 0 & 1 & 0 \\
U12 & 1 & 0 & 0 & 1 & 0 & 0 & 1 & 0 \\
U13 & 0 & 1 & 0 & 0 & 0 & 1 & 0 & 1 \\
U14 & 0 & 1 & 0 & 1 & 0 & 1 & 0 & 1 \\
U15 & 0 & 0 & 1 & 0 & 1 & 0 & 1 & 0 \\
U16 & 0 & 0 & 0 & 1 & 1 & 1 & 1 & 0 \\
\bottomrule
\end{tabular}

\medskip
\begin{minipage}{0.9\textwidth}
\footnotesize\raggedright\textit{Note.} Memberships were coded by the authors from the interviews, using the rules in Section 3.3. Synthetic data.
\end{minipage}
\end{table}

\section{Findings}

\subsection{Necessity}

No condition is necessary for reliable handoffs. We tested each of the seven conditions and its negation against a necessity consistency threshold of 0.90. None passed: no condition, and no negated condition, holds for more than six of the eight reliable units. We therefore make no claim that any condition is necessary.

\subsection{Sufficiency}

Six of the eight reliable units share all seven codes with a unit outside the outcome set. Each of these rows therefore holds one reliable unit and one other unit and has a raw consistency of 0.50, below our cutoff, and we coded these contradictory rows 0. This is why the two configurations in Table~\ref{tab:fit} cover two of the eight reliable units. The two rows coded 1 each hold one unit, and both units are reliable.

\begin{table}[htbp]
\centering
\caption{Configurations sufficient for reliable handoffs (conservative solution)}
\label{tab:fit}
\begin{tabular}{lcc}
\toprule
Condition & Configuration 1 & Configuration 2 \\
\midrule
Protocol & $\bullet$ & $\otimes$ \\
Champion & $\otimes$ & $\bullet$ \\
Huddles & $\bullet$ & $\bullet$ \\
Staffing & $\otimes$ & $\bullet$ \\
Tenure & $\otimes$ & $\otimes$ \\
Digital & $\bullet$ & $\otimes$ \\
Leader & $\otimes$ & $\otimes$ \\
\midrule
Raw consistency & 1.00 & 1.00 \\
Raw coverage & 0.12 & 0.12 \\
Unique coverage & 0.12 & 0.12 \\
\midrule
Solution consistency & \multicolumn{2}{c}{1.00} \\
Solution coverage & \multicolumn{2}{c}{0.25} \\
\bottomrule
\end{tabular}

\medskip
\begin{minipage}{0.9\textwidth}
\footnotesize\raggedright\textit{Note.} $\bullet$ = condition present; $\otimes$ = condition absent. Configuration 1 is Protocol*$\sim$Champion*Huddles*$\sim$Staffing*$\sim$Tenure*Digital*$\sim$Leader; configuration 2 is $\sim$Protocol*Champion*Huddles*Staffing*$\sim$Tenure*$\sim$Digital*$\sim$Leader. Figures are rounded to two decimals, half to even. Synthetic data.
\end{minipage}
\end{table}

Table~\ref{tab:fit} reports two configurations, each with a raw consistency of 1.00. We call the first configuration the protocol-led unit. A written protocol is in use, there is no named champion, huddles are held, and handoffs are recorded in the electronic tool. We call the second configuration the champion-led unit. A champion is named, there is no protocol in use, huddles are held, and staffing is adequate. The two configurations share no unit, since one requires a protocol and the other requires its absence. Each covers one of the eight reliable units, a coverage of one eighth that Table~\ref{tab:fit} prints as 0.12 after rounding. Together they cover two of the eight reliable units, which the table prints as 0.25.

Where a protocol carried the handoff, informants described it as a script that did not depend on who was on shift. A charge nurse said: ``The sheet does the remembering. It doesn't matter if I'm on or a traveler is on, we go down the same lines in the same order'' (charge nurse, interview 31). A resident made the same point from the other side of the handoff: ``I get the same five things every morning, in the same order, so I notice when one is missing'' (resident physician, interview 47).

Where a champion carried the handoff, informants described a person who watched and corrected. ``If you skip the pending labs, she'll be at the station before you've left, asking what you forgot'' (staff nurse, interview 12). A physician described the same role: ``Everyone knows who to ask, and everyone knows she'll ask them'' (attending physician, interview 58).

Neither configuration contains both a protocol and a champion. Huddles are present in both. Long tenure and a manager at shift change are absent in both, so neither configuration rests on a stable, closely supervised workforce.

\section{Discussion}

We set out to ask whether reliable handoffs require both a rule and a person. In this hospital they do not. Two configurations are sufficient for reliable handoffs, the protocol-led unit and the champion-led unit, and each contains one of the two and the absence of the other. The finding is about sufficiency. It does not say that protocols or champions matter in every unit, and the two configurations cover only two of the eight reliable units.

The other six reliable units sit in contradictory rows. Each shares its codes with a unit that did not sustain reliable handoffs, so the seven conditions do not separate those units from their counterparts. Something we did not code distinguishes them. Informants in those units spoke most often about particular physicians and about the mix of patients, and a later study could code those conditions.

\subsection{Robustness}

We checked three choices. First, any raw consistency cutoff above 0.50 and up to 1.00 codes the same rows as our cutoff of 0.80, because every observed row has a consistency of 1.00, 0.50, or 0. Second, a frequency cutoff of two units would leave no row coded 1, since each configuration rests on a single unit. We kept the cutoff at one unit and report this here so readers can weigh it. Third, we recoded each of the two borderline codes, one at a time. Neither change moved a unit into or out of either configuration, and both configurations survived unchanged.

\subsection{Limitations}

We studied one hospital over one year, so the configurations describe these units and are not a claim about handoffs elsewhere. We analyzed the presence of reliable handoffs only. Our interviews asked what sustained reliable handoffs, and informants in the other units described many separate failures that our seven conditions were not designed to capture, so we make no claim about the absence of the outcome. Our coding rests on interviews, and a unit's code is only as good as the accounts its informants gave us. We tried to guard against this by requiring agreement across roles for the outcome, the champion, and staffing.

\begin{thebibliography}{9}
\setlength{\itemsep}{0pt}
\bibitem[Fiss(2011)]{fiss2011} Fiss, P. C. (2011). Building better causal theories: A fuzzy set approach to typologies in organization research. \textit{Academy of Management Journal}, 54(2), 393--420.
\bibitem[Greckhamer et al.(2018)]{greckhamer2018} Greckhamer, T., Furnari, S., Fiss, P. C., and Aguilera, R. V. (2018). Studying configurations with qualitative comparative analysis: Best practices in strategy and organization research. \textit{Strategic Organization}, 16(4), 482--495.
\bibitem[Ragin(1987)]{ragin1987} Ragin, C. C. (1987). \textit{The comparative method: Moving beyond qualitative and quantitative strategies}. University of California Press.
\bibitem[Ragin(2008)]{ragin2008} Ragin, C. C. (2008). \textit{Redesigning social inquiry: Fuzzy sets and beyond}. University of Chicago Press.
\bibitem[Schneider and Wagemann(2012)]{schneider2012} Schneider, C. Q., and Wagemann, C. (2012). \textit{Set-theoretic methods for the social sciences: A guide to qualitative comparative analysis}. Cambridge University Press.
\end{thebibliography}

\end{document}
